\section{Metodología}

\subsection{Caracterización de la probeta}

Se midió la profundidad de la probeta con una regla y su diámetro interno con un calibrador Vernier. La temperatura ambiente se registró con un termómetro. Estas mediciones se utilizaron para obtener la longitud efectiva del tubo y calcular la velocidad de referencia del sonido en aire mediante la Ec.~\eqref{eq:velocidad_temperatura}.

% TODO para el equipo: completar la resolución/incertidumbre de la regla, el
% calibrador Vernier y el termómetro; indicar cuántas mediciones se repitieron.

\subsection{Identificación de las frecuencias de resonancia}

Se utilizó un teléfono celular con una aplicación generadora de frecuencias y otra para analizar el espectro. El teléfono se colocó sobre la abertura de la probeta, con el micrófono dentro del tubo, como se muestra en la Fig.~\ref{fig:montaje_frecuencias}. Se modificó la frecuencia de excitación y se localizaron en el espectro los máximos asociados con las resonancias; las frecuencias correspondientes se registraron para identificar los modos y sus números de armónico. La Fig.~\ref{fig:esquema_montaje} presenta un esquema del montaje, sin incluir conexiones adicionales.

% TODO para el equipo: confirmar con el registro original la configuración aplicada.
% La guía indica barridos de 150 a 2150 Hz durante 10 s, repetidos dos veces,
% mientras que la Tabla 3 consigna 100 a 2300 Hz durante 20 s.

\begin{figure}[!t]
\centering
\begin{minipage}[t]{0.48\columnwidth}
\centering
\includegraphics[width=\linewidth,height=0.16\textheight,keepaspectratio]{probetaCelular.jpeg}
\caption{Montaje con teléfono para identificar las frecuencias de resonancia.}
\label{fig:montaje_frecuencias}
\end{minipage}\hfill
\begin{minipage}[t]{0.48\columnwidth}
\centering
\includegraphics[width=\linewidth,height=0.16\textheight,keepaspectratio]{probetaagua.jpeg}
\caption{Probeta con agua y columna de aire.}
\label{fig:columna_aire}
\end{minipage}
\end{figure}

\begin{figure}[!t]
\centering
\begin{minipage}[c]{0.54\columnwidth}
\centering
\shorthandoff{<>}
\begin{tikzpicture}[
  x=1cm,y=0.52cm,
  font=\scriptsize,
  wall/.style={line width=1pt},
  dimension/.style={{Stealth[reversed]}-{Stealth}, thin},
  label/.style={align=center}
]
  % Probeta abierta arriba; el agua representa el cierre acústico inferior.
  \draw[wall] (0,0) -- (0,5.4);
  \draw[wall] (1.5,0) -- (1.5,5.4);
  \draw[wall] (0,0) -- (1.5,0);
  \fill[blue!20] (0.02,0.03) rectangle (1.48,1.25);
  \draw[blue!60!black,thick] (0.02,1.25) -- (1.48,1.25);
  \node[label] at (0.75,0.65) {agua};
  \node[label] at (0.75,3.2) {columna de aire};
  \draw[dimension] (1.8,1.3) -- (1.8,5.3);
  \node[right,label] at (1.88,3.3) {$L'$};
  \node[above,label] at (0.75,5.45) {Extremo abierto};
  \node[below,label] at (0.75,-0.08) {Cierre acústico};
  \draw[wall,rounded corners=2pt] (-1.8,3.8) rectangle (-0.15,5.2);
  \node[label] at (-0.975,4.5) {Teléfono\\generador y\\analizador};
  \draw[->] (-0.15,4.8) -- (0.12,5.35);
\end{tikzpicture}
\shorthandon{<>}
\caption{Esquema del montaje experimental.}
\label{fig:esquema_montaje}
\end{minipage}\hfill
\begin{minipage}[c]{0.43\columnwidth}
\centering
\includegraphics[width=\linewidth,height=0.16\textheight,keepaspectratio]{diapasonprobeta.jpeg}
\caption{Comprobación de resonancia con diapasón.}
\label{fig:resonancia_diapason}
\end{minipage}
\end{figure}

Una vez identificados los modos, se relacionaron las frecuencias con sus números de armónico. El ajuste lineal de $f_n$ respecto de $n$ permitió estimar la frecuencia fundamental y, mediante el modelo de tubo cerrado, obtener una estimación experimental de la velocidad del sonido.

\subsection{Variación de la columna de aire}

Se agregó agua progresivamente a la probeta para reducir la longitud de la columna de aire, como se ilustra en la Fig.~\ref{fig:columna_aire}. Para cada nivel se buscó y registró la frecuencia fundamental de resonancia. La longitud restante se relacionó con el volumen agregado mediante la Ec.~\eqref{eq:longitud_aire_agua}; posteriormente, la dependencia de la frecuencia respecto de la longitud inversa se utilizó para estimar la velocidad del sonido a partir de la pendiente del ajuste.

% TODO para el equipo: confirmar los volúmenes realmente utilizados. La Tabla 6
% de la guía indica 0, 75, 150, 225, 300, 375 y 450 cm^3; el procedimiento
% también menciona una adición inicial de 50 mL. No se atribuyen aquí esos
% valores a las mediciones realizadas hasta contrastarlos con los registros.

\subsection{Comprobación de resonancia con diapasones}

Se utilizaron diapasones de frecuencias conocidas. Para cada frecuencia se calculó el volumen de agua requerido con la Ec.~\eqref{eq:volumen_resonancia}; luego se ajustó el nivel de agua y se comprobó experimentalmente la resonancia, según se muestra en la Fig.~\ref{fig:resonancia_diapason}.

% TODO para el equipo: confirmar las frecuencias de los diapasones efectivamente
% usados. La guía enumera 256, 320, 384 y 480 Hz.
